\section{Bugs and Incidents: Fast Lanes and Learning}
Bugs and incidents need speed and structure. Treat them as a separate flow with severity rules, fast lanes, and a postmortem pattern.

\subsection{Jira setup steps (Bugs/Incidents)}
\begin{enumerate}
  \item Create a Bug issue type and add Severity (S1–S4) as a field or label.
  \item Configure a Bug workflow: Triage \textrightarrow{} Fix In Progress \textrightarrow{} Verify \textrightarrow{} Done, with Won’t Fix.
  \item Create a fast-lane swimlane for Severity S1/S2.
  \item Add a postmortem issue type or Task template linked to the Bug.
  \item Automate creation of a postmortem Task when Severity is S1.
\end{enumerate}

\subsection{Severity model}
S1: customer outage or data loss. S2: major degradation. S3: minor defect. S4: cosmetic. Severity determines SLA and escalation path.

Sample JQL filters:
\begin{itemize}
  \item Active S1/S2: \texttt{project = PMH AND issuetype = Bug AND Severity in (S1, S2) AND status != Done}
  \item Bugs past SLA: \texttt{project = PMH AND issuetype = Bug AND status != Done AND updated <= -2d}
\end{itemize}

\subsection{Fast lanes}
Fast lanes are swimlanes or quick filters that surface critical issues. For S1/S2, bypass standard WIP limits but require EM or TPM approval to enter.

\subsection{Postmortem issue pattern}
For every S1, create a linked Task called "Postmortem: <incident>" with owners and due dates. Track root cause, fixes, and follow-up items.

Sample JQL filters:
\begin{itemize}
  \item Missing postmortem: \texttt{project = PMH AND issuetype = Bug AND Severity = S1 AND issueLinkType is EMPTY}
\end{itemize}

\subsection{Case study insert}
BG-124 is marked S1 due to failed upgrade emails. The bug moves through the fast lane and spawns a postmortem task linked to EPIC-5.






